\section{Inicializace a stuktura projektu}
Pro urychlení počáteční fáze vývoje a zajištění moderní architektury prezentační vrstvy byl využit Vue starter kit, který vytvořili samotní tvůrci Laravel frameworku. Toto řešení poskytuje základní kostru aplikace včetně autentizační logiky a propojení backendu s frontendovým frameworkem Vue.js. Vizuální stránka uživatelského rozhraní je zajišťována pomocí knihovny \mbox{shadcn-vue}, která nabízí sadu přizpůsobitelných komponent. Pro jednotný vizuální styl ovládacích prvků je dále využívána sada open source ikon Lucide.

Projekt dodržuje konvence Laravel frameworku pro adresářovou strukturu, což zajišťuje přehlednost a snadnou udržitelnost kódu. Klíčové součásti implementace jsou znázorněny na následujícím schématu:

\dirtree{%
  .1 .
  .1 app\DTcomment{jádro aplikace obsahující aplikační logiku}.
    .2 Http .
      .3 Controllers\DTcomment{kontrolery zpracovávající příchozí požadavky}.
      .3 Middleware\DTcomment{filtry HTTP požadavků (např. autorizace)}.
      .3 Requests\DTcomment{typy požadavků s validačními pravidly}.
    .2 Models\DTcomment{definice modelů reprezentující datové entity}.
    .2 Services.
      .3 GitProviders\DTcomment{služby pro synchronizaci dat s VCS}.
  .1 config\DTcomment{konfigurační soubory aplikace}.
  .1 database\DTcomment{databázové migrace, seedery pro testovací data}.
  .1 resources\DTcomment{zdrojové kódy prezentační vrstvy (CSS a JS)}.
  .1 routes\DTcomment{směrovací pravidla}.
  .1 .env\DTcomment{proměnné prostředí}.
  .1 composer.json\DTcomment{soubor se seznamem PHP závislostí}.
  .1 package.json\DTcomment{soubor se seznamem JavaScriptových závislostí}.
}
